\documentclass[reqno, 12pt]{amsart}

\usepackage{amssymb}
\usepackage{amsthm}
\usepackage{amsmath}
\usepackage{amsxtra}
\usepackage{mathrsfs}
\usepackage{comment}
\usepackage{graphicx}
\usepackage{placeins}
\usepackage{multicol}
\usepackage{bbm}
\usepackage{stmaryrd}
%\usepackage{enumerate}
\usepackage[inline, shortlabels]{enumitem}
\usepackage{mathdots}
\usepackage{colonequals}
\usepackage{hyperref}


\usepackage{calc}
\usepackage{tikz}
\usetikzlibrary{arrows,calc,automata,shadows,backgrounds,positioning,intersections,fadings,decorations.pathreplacing,shapes,matrix}
\usepackage{tikz-cd}

\usepackage{fullpage}

\newtheorem{thm}{Theorem}[section]
\newtheorem{lem}[thm]{Lemma}
\newtheorem{cor}[thm]{Corollary}
\newtheorem{conj}[thm]{Conjecture}
\newtheorem{prop}[thm]{Proposition}
\theoremstyle{definition}  
\newtheorem{defn} [thm] {Definition} 
\newtheorem{claim}[thm]{Claim}
\newtheorem{example} [thm] {Example}
\newtheorem{rem} [thm] {Remark}

\newenvironment{problab}[1]
{\noindent\textbf{Problem #1}.}
{\vskip 6pt}

\newenvironment{exolab}[1]
{\noindent\textbf{Exercise #1}.}
{\vskip 6pt}

\theoremstyle{remark}
\newtheorem*{soln}{Solution}

\newcommand{\C}{\mathbb C}
\renewcommand{\H}{\mathbb H}
\newcommand{\D}{\mathbb D}
\newcommand{\F}{\mathbb F}
\newcommand{\Q}{\mathbb Q}
\newcommand{\R}{\mathbb R}
\newcommand{\Z}{\mathbb Z}
\newcommand{\T}{\mathbb T}
\renewcommand{\P}{\mathbb P}
\renewcommand{\O}{\mathcal O}
\newcommand{\m}{\mathfrak m}
\newcommand{\A}{\mathbb A}
\renewcommand{\SS}{\mathbb S}

\newcommand{\B}{\mathcal B}
\newcommand{\E}{\mathcal E}

\newcommand{\SSpan}{\text{span}}

\DeclareMathOperator{\lcm}{lcm}
\DeclareMathOperator{\img}{img}
\DeclareMathOperator{\Ann}{Ann}
\DeclareMathOperator{\Tor}{Tor}
\DeclareMathOperator{\tr}{tr}
\DeclareMathOperator{\Hom}{Hom}
\DeclareMathOperator{\End}{End}
\DeclareMathOperator{\Aut}{Aut}
\DeclareMathOperator{\Gal}{Gal}
\DeclareMathOperator{\sgn}{sgn}
\DeclareMathOperator{\ord}{ord}
\DeclareMathOperator{\ad}{ad}
\DeclareMathOperator{\coker}{coker}
\DeclareMathOperator{\rank}{rank}

\DeclareMathOperator{\id}{id}

\usepackage{mathpazo}

\everymath{\displaystyle}

\tikzset{commutative diagrams/.cd, arrow style = tikz, diagrams = {>=latex}}

\newenvironment{amatrix}[1]{%
  \left(\begin{array}{@{}*{#1}{c}|c@{}}
}{%
  \end{array}\right)
}

\usepackage{abstract}
\renewcommand{\abstractname}{}    % clear the title
\renewcommand{\absnamepos}{empty} % originally center
%\setlength{\abstitleskip}

\begin{document}
%\renewcommand{\abstractname}{\vspace{-\baselineskip}}

\title{18.700 Problem Set 3}
%\author{Évariste Galois}
\dedicatory{Due Wednesday, September 25 at 11:59 pm on \href{https://canvas.mit.edu/courses/27315}{Canvas}}
%\thanks{}
%\contrib{}
%\address{}

\maketitle

\begin{abstract}
  \noindent Collaborated with:\\
  Sources used: 
\end{abstract}
\vskip 0.5cm

\begin{problab}{1} (2A \#5) (5 points)
  \begin{enumerate}[(a)]
    \item 
      Find a number $t \in \R$ such that
      \begin{equation*}
        (3,1,4),\ (2,-3,5),\ (5,9,t)
      \end{equation*}
      is linearly dependent in $\R^3$.
    \item
      For the value of $t$ found in the previous part, determine the smallest $k \in \{1,2,3\}$ such that $v_k \in \SSpan(v_1, \ldots, v_{k-1})$. Express the $v_k$ as a linear combination of $v_1, \ldots, v_{k-1}$.
  \end{enumerate}
\end{problab}

\begin{problab}{2} (2A \#7) (5 points)
  \begin{enumerate}[(a)]
    \item 
      Show that if we consider $\C$ as vector space over $\R$, the list $1+i, 1-i$ is linearly independent.
    \item
      Show that if we consider $\C$ as vector space over $\C$, the list $1+i, 1-i$ is linearly dependent.
  \end{enumerate}
\end{problab}

\begin{problab}{3} (2A \#13) (6 points)
  Suppose that $v_1, \ldots, v_m \in V$ are linearly independent and $w \in V$. Show that $v_1, \ldots, v_m, w$ are linearly independent iff $w \notin \SSpan(v_1, \ldots, v_m)$.
\end{problab}

\begin{problab}{4} (1C \#24) (10 points) % p. 26
  Let $f: \R \to \R$ be a function.
  \begin{itemize}
    \item $f$ is called \textit{even} if $f(-x) = f(x)$ for all $x \in \R$.
    \item $f$ is called \textit{odd} if $f(-x) = -f(x)$ for all $x \in \R$.
  \end{itemize}
  Let
  \begin{align*}
    V_e &\colonequals \{f: \R \to \R \mid \text{$f$ is even}\}, \ \text{and}\\
    V_o &\colonequals \{f: \R \to \R \mid \text{$f$ is odd}\} \, .
  \end{align*}
  \begin{enumerate}[(a)]
    \item 
      Show that $V_e$ and $V_o$ are subspaces of $\R^{\R}$.
    \item
      Show that $\R^{\R} = V_e \oplus V_o$.
  \end{enumerate}
\end{problab}

\begin{problab}{5} (11 points)
  Consider the following subspaces of $\F^3$:
  \begin{align*}
    V_1 &\colonequals \{(x,y,0) \in \F^3 \mid x,y \in \F\}\\
    V_2 &\colonequals \{(x,0,z) \in \F^3 \mid x,z \in \F\}\\
    V_3 &\colonequals \{(0,y,z) \in \F^3 \mid y,z \in \F\} \, .
  \end{align*}
  \begin{enumerate}[(a)]
    \item Compute $V_1 \cap V_2 \cap V_3$.
    \item Is the sum $V_1 + V_2 + V_3$ direct? Why or why not?
    \item Prove the following generalized criterion for a sum to be direct. Let $V$ be a vector space and $V_1, \ldots, V_m$ be subspaces of $V$. Then the sum $V_1 + \cdots + V_m$ is direct iff
      \begin{equation*}
        V_j \cap \left( \sum_{i \neq j} V_i \right) = V_j \cap \left(V_1 + \cdots V_{j-1} + V_{j+1} + \cdots + V_m \right) = \{0\}
      \end{equation*}
      for each $j = 1, \ldots, m$.
  \end{enumerate}
\end{problab}

\end{document} 	
